\documentclass[10pt,a4paper]{amsart}
\usepackage[margin=19mm]{geometry}
\usepackage{amsmath,amssymb,mathtools}
\usepackage[scr=rsfs,cal=euler]{mathalfa}
\usepackage{xfrac}
\usepackage[unicode,hidelinks,bookmarks=false]{hyperref}
\newcommand{\ie}{i.e.\ }
\newcommand{\cf}{cf.\ }
\newcommand{\resp}{respectively\ }
\newcommand{\Subsection}{\subsection}
\providecommand{\nobreakdash}{\nobreak\hbox{-}}
\setlength{\emergencystretch}{2em}
\multlinegap0pt
\usepackage{amssymb}
\usepackage{enumerate}
\usepackage[shortlabels]{enumitem}
\newtheorem{lemma}{Lemma}
\newtheorem{theorem}{Theorem}
\renewcommand{\leq}{\leqslant}
\title{Some Series related to Extended Riemann Hypothesis for Dedekind zeta functions}
\author{Muhammad Atif Zaheer}
\date{Source: arXiv:2512.19761v1 (CC BY 4.0)}
\begin{document}
\maketitle

\noindent\emph{Source and licence.} Some Series related to Extended Riemann Hypothesis for Dedekind zeta functions. Version arXiv:2512.19761v1 (CC BY 4.0), distributed by arXiv under the Creative Commons Attribution 4.0 International licence, http://creativecommons.org/licenses/by/4.0/, as recorded in the arXiv licence metadata for this exact version and shown on the abstract page https://arxiv.org/abs/2512.19761v1. Full licence notice: \texttt{licenses/arxiv-2512.19761-CC-BY-4.0.txt}.\par\smallskip

\noindent\textit{Editorial scope.} Counted unit: the article's theorem result (source0.tex lines 321--327). The article's complete original proof of it is delivered with it (source0.tex lines 329--369, one \texttt{proof} environment of 41 lines). No mathematics is written, repaired or re-derived: the statement and the proof are the article's own lines, in the article's own notation, and no retained line is substituted or re-flowed.  Retained uncounted as the source-local context the proof needs: lines 291--293; lines 298--300.  Nothing was omitted from any retained span, so the package carries no omission record.  No retained line contains a citation, so the wrapper carries no bibliography and no bibliography entry.  No retained line contains a figure environment, an image-inclusion command or drawing code, so no image is delivered and the package declares no asset dependency.  Editorial formatting only, in addition: an AMS \texttt{amsart} preamble that restates the article's own declarations and macros used by the retained lines, namely the environments \texttt{lemma}, \texttt{theorem} on separate counters, together with the standard packages those lines need.  The retained environments print in this document as Lemma 1, Theorem 1 in order of appearance, whereas the article prints the same items with the numbering of its own full document.  The complete native source of the article as distributed in the arXiv e-print for this exact version is bundled unchanged as \texttt{sources/191505/source0.tex}.
\begin{lemma} \label{vanishing of X_K at 1/2} 
For $k$ even, $(\Xi_K'/\Xi_K)^{(k)}(1/2) = 0$.
\end{lemma}

\begin{equation} \label{hadamard product}
\Xi_K(s) = e^{A + Bs} \prod_{\rho_K \neq 1/2} \left( 1 - \frac{s}{\rho_K} \right) e^{s/\rho_K},
\end{equation}

\begin{theorem}
Let $K$ be a number field and let $k > 1$ be an odd positive integer. If ERH holds, then 
\begin{equation} \label{main result}
\frac{(-1)^{(k - 1)/2}}{k!} \left( \frac{\Xi_K'}{\Xi_K} \right)^{(k)} \left(\frac{1}{2} \right) = \sum_{\rho_K \neq 1/2} \frac{1}{|1/2 - \rho_K|^{k + 1}},
\end{equation}
where the sum runs over all nontrivial zeros $\rho_K$ of $\zeta_K(s)$ distinct from $1/2$. Moreover, if the above equality holds, then $\beta_K = 1/2$ for $|\gamma_K| > (1/2)\cot (2 \pi/(k + 1))$.
\end{theorem}

\begin{proof}
Taking the logarithmic derivative of the Hadamard product \eqref{hadamard product} we obtain
\[
\frac{\Xi_K'}{\Xi_K}(s) = B + \sum_{\rho_K \neq 1/2} \left( \frac{1}{s - \rho_K} + \frac{1}{\rho_K} \right).
\]
Because the sum converges uniformly on compact subsets not containing the nontrivial zeros of $\zeta_K(s)$ we can differentiate term by term and deduce that
\[
\left( \frac{\Xi_K'}{\Xi_K} \right)'(s) = - \sum_{\rho_K \neq 1/2} \frac{1}{(s - \rho_K)^2}.
\]
By repeated differentiation we end up with the identity
\[
\frac{(-1)^k}{k!}\left( \frac{\Xi_K'}{\Xi_K} \right)^{(k)}(s) = \sum_{\rho_K \neq 1/2} \frac{1}{(s - \rho_K)^{k + 1}}
\]
which holds for any positive integer $k$. For $k$ even it immediately follows from Lemma \ref{vanishing of X_K at 1/2} that
\[
\sum_{\rho_K \neq 1/2} \frac{1}{(1/2 - \rho_K)^{k + 1}} = 0.
\]
But this is well-known and follows simply due to the symmetry of zeros $\rho_K$ under conjugation and reflection across the critical line. We now assume that $k$ is odd. By noting that $(\Xi_K'/\Xi_K)^{(k)} (1/2)$ is real we obtain
\begin{align*}
\frac{(-1)^k}{k!}\left( \frac{\Xi_K'}{\Xi_K} \right)^{(k)} \left( \frac{1}{2} \right) &= \frac{1}{2}\sum_{\rho_K \neq 1/2} \frac{1}{(1/2 - \rho_K)^{k + 1}} + \frac{1}{2} \sum_{\rho_K \neq 1/2} \frac{1}{(1/2 - \overline{\rho_K})^{k + 1}}  \\
&= \frac{1}{2} \sum_{\rho_K \neq 1/2} \frac{(1/2 - \rho_K)^{k + 1} + (1/2 - \overline{\rho_K})^{k + 1}}{|1/2 - \rho_K|^{2(k + 1)}} \\
&= \sum_{\rho_K \neq 1/2} \frac{\Re(1/2 - \rho_K)^{k + 1}}{|1/2 - \rho_K|^{2(k + 1)}}.
\end{align*}
If ERH holds, then $\Re(1/2 - \rho_K)^{k + 1} = (-1)^{(k + 1)/2} |1/2 - \rho_K|^{k + 1}$ and so we get \eqref{main result}.

Now suppose that \eqref{main result} holds. Then we have 
\begin{equation} \label{second main equation}
\sum_{\rho_K \neq 1/2} \frac{\Re(1/2 - \rho_K)^{k + 1} - (-1)^{(k + 1)/2}|1/2 - \rho_K|^{k + 1}}{|1/2 - \rho_K|^{2(k + 1)}} = 0.
\end{equation}
If $k \equiv 1 \pmod{4}$, then the numerator in the summand is always nonnegative and so we have $\Re(1/2 - \rho_K)^{k + 1} = - |1/2 - \rho_K|^{k + 1}$ for every zero $\rho_K \neq 1/2$. Let $1/2 - \rho_K = |1/2 - \rho_K| e({ \phi_{\rho_K}})$, where we use $e(x)$ to denote $e^{ix}$ and $\phi_{\rho_K}$ is the principal argument of $1/2 - \rho_K$. Then it follows that
\[
\Re \left( \left|1/2 - \rho_K \right|^{k + 1} e({(k + 1) \phi_{\rho_K}}) \right)= - \left|1/2 - \rho_K \right|^{k + 1}
\]
for every zero $\rho_K \neq 1/2$, which in turn implies that $\cos((k + 1) \phi_{\rho_K}) = -1$ and thus $\phi_{\rho_K} =((2n + 1) \pi)/(k + 1)$, where $0 \leq n \leq  k$. Because $-1/2 < 1/2 - \beta_K < 1/2$, it follows that if $|\gamma_K|> (1/2) \cot (2 \pi/(k + 1))$, then $\tan^{-1}(2|\gamma_K|) > \pi/2 - 2\pi/(k + 1)$. This yields $|\phi_{\rho_K} - \pi/2| < 2 \pi/(k + 1)$  or $|\phi_{\rho_K} - 3 \pi/2| < 2 \pi/(k + 1)$. Consequently we must have either $\phi_{\rho_K} = \pi/2$ or $3 \pi/2$  and hence $1/2- \beta_K =  \Re(1/2  -\rho_K) = |1/2 - \rho_K|\cos(\phi_{\rho_K}) = 0$, i.e.,  $\beta_K = 1/2$.

If $k \equiv 3 \pmod{4}$. Then numerator in each summand in the sum \eqref{second main equation} is nonpositive for every zero $\rho_K \neq 1/2$. Thus $\Re (1/2 - \rho_K)^{k + 1} = |1/2 - \rho_K|^{k + 1}$ and so
\[
\Re \left(\left| 1/2 - \rho_K \right|^{k + 1} e({(k + 1) \phi_{\rho_K}}) \right)=  \left| 1/2 - \rho_K \right|^{k + 1}
\]
for every zero $\rho_K \neq 1/2$, which leads to $\cos((k + 1) \phi_{\rho_K}) = 1$ and so $\phi_{\rho_K} = 2\pi n/(k + 1)$, where $0 \leq n \leq k$. Again it follows that if $|\gamma_K| > (1/2) \cot (2 \pi/(k + 1))$, then we must have $\phi_{\rho_K} =  \pi/2$ or $3 \pi/2$. Hence $1/2 - \beta_K = \Re(1/2 - \rho_K) = |1/2 - \rho_K| \cos (\phi_{\rho_K}) = 0$, i.e., $\beta_K = 1/2$. This completes the proof.
\end{proof}

\end{document}
